% Supplementary material: moved from Sections 10.1, 10.2, 10.4 and 10.5 of the main text.

\section{Model-Level Challenges in Detail}
\label{supp:challenges}

\subsection{Hallucination Mitigation in Small Models}
\label{subsec:challenges_hallucination}

Hallucination, the confident generation of factually incorrect content, is an intrinsic failure mode of autoregressive language models, and it is more pronounced in smaller models \citep{zhang2023siren}. The core issue is information capacity: a 3B model's parameters can encode far fewer facts than a 175B model, and when queried about knowledge at the boundary of its training distribution, small models are more likely to confabulate plausible-sounding but incorrect responses.

Mitigation strategies include: (1)~\emph{\gls{rag}} \citep{lewis2020rag}, which anchors \gls{slm} generation to retrieved documents rather than parametric memory; (2)~\emph{uncertainty quantification} through sampling and ensemble methods; (3)~\emph{RLHF with truthfulness rewards} to penalize confident incorrect responses; and (4)~\emph{calibration post-processing} to align model confidence to actual accuracy.

\subsection{Reasoning Limitations Under Parameter Constraints}
\label{subsec:challenges_reasoning}

Language-modeling loss falls as a power law in parameter count \citep{kaplan2020scaling}, and in the reported results reviewed here the gap between small and large models is widest on multi-step reasoning, particularly in mathematics, formal logic, and scientific inference (\evD). At the 1B--7B scale, \glspl{slm} struggle with problems that require holding many intermediate states simultaneously or with deep deductive chains.

Emerging directions include: (1)~\emph{Test-time scaling}: allocating additional inference compute (more generation tokens, repeated sampling, tree-of-thought search) to substitute for model capacity; the DeepSeek-R1 authors \citep{deepseek2025r1} report that distilling the long chains of thought of their large reasoning model into 1.5B--70B students transfers much of its mathematical-reasoning ability (\evC). (2)~\emph{Neuro-symbolic integration}: augmenting \glspl{slm} with symbolic solvers (SAT solvers, theorem provers, calculators) to handle formal reasoning externally. (3)~\emph{Speculative decoding} \citep{leviathan2023speculative,chen2023specsampling} using a small draft \gls{slm} to propose tokens verified by a larger model, which speeds up the larger model rather than substituting for it (Section~\ref{subsubsec:agentic_routing}).

\subsection{Hardware Co-Design and NPU Optimization}
\label{subsec:challenges_hardware}

The efficiency of \gls{slm} inference depends as much on hardware as on model architecture. Modern consumer hardware increasingly includes dedicated AI accelerators: Apple ANE, Qualcomm Hexagon \gls{npu} \citep{qualcomm2024ondevice}, Samsung Exynos \gls{npu}, and Intel \gls{npu}. These accelerators are optimized for specific operation types (INT4/INT8 matrix multiplication) and data layouts.

A critical open problem is \emph{hardware-model co-design}: developing \gls{slm} architectures whose attention patterns, FFN widths, and layer counts align with the data bus widths and tile sizes of target \glspl{npu}. Future \glspl{slm} trained with hardware constraints as first-class objectives, rather than retrofitting quantization after the fact, will likely yield significant efficiency improvements.

\subsection{Ethical Considerations: Bias Amplification in Compressed Models}
\label{subsec:challenges_ethics}

Model compression introduces an underexplored ethical risk: \emph{bias amplification}. Pruning and quantization may disproportionately affect model performance on underrepresented demographic groups, dialects, and languages, as the model's remaining capacity may preferentially serve the majority distributions present in training data \citep{blodgett2020language}. \citet{kurz2025pruning_vs_quant} find that quantization preserves multilingual perplexity better than pruning, suggesting that pruning disproportionately degrades non-English performance.

Research priorities include: (1)~\emph{fairness-aware compression} that includes demographic parity metrics as part of the compression objective; (2)~\emph{representation auditing} of compressed models across language, dialect, gender, and race; and (3)~\emph{minimum performance floors} for safety-critical applications, ensuring compressed models meet baseline accuracy thresholds for all user populations.

\subsection{Sustainable AI: SLMs as a Green AI Alternative}
\label{subsec:challenges_sustainability}

The energy consumption of training and serving frontier \glspl{llm} has become a significant environmental and economic concern \citep{strubell2019energy,patterson2022carbon}. Training GPT-3 emitted an estimated 552 metric tons of CO$_2$ equivalent \citep{patterson2021carbon}, and training a model two orders of magnitude smaller from scratch needs correspondingly less compute. The defensible claim about inference is conditional: \emph{for a task a small model completes correctly, on hardware it keeps busy, serving it with the small model uses less inference energy} (\evD). Whether a shift to small models reduces total system energy depends on teacher-training energy, the compression pipeline, device embodied carbon, regional grid intensity, update distribution, utilization, and the rebound effect, and no published accounting answers it. \suppsec~\ref{supp:energy} sets out the accounting boundary.

\subsection{Democratization and Sustainability}
\label{supp:democratization}

Running a 3B model with strong reported grade-school-mathematics accuracy on an
inexpensive single-board computer is a qualitative change in access.
Organizations facing cloud-API cost, data-sovereignty rules, or unreliable
connectivity can run capable \gls{nlp} locally with no per-query charge and,
subject to the scope conditions of Section~\ref{subsec:deployment_privacy_scope},
without transferring prompt content to a processor outside their control. The
environmental side of the argument is conditional in the way set out above and in
\suppsec~\ref{subsubsec:challenges_env_scope}: we make a per-task inference claim
and no aggregate one.

